
\documentclass{article}
\usepackage{colortbl}
\usepackage{makecell}
\usepackage{multirow}
\usepackage{supertabular}

\begin{document}

\newcounter{utterance}

\centering \large Interaction Transcript for game `cladder', experiment `full\_v1.5\_default', episode 9771 with llama{-}3.1{-}8b{-}instruct.
\vspace{24pt}

{ \footnotesize  \setcounter{utterance}{1}
\setlength{\tabcolsep}{0pt}
\begin{supertabular}{c@{$\;$}|p{.15\linewidth}@{}p{.15\linewidth}p{.15\linewidth}p{.15\linewidth}p{.15\linewidth}p{.15\linewidth}}
    \# & \multicolumn{2}{c}{Player} && \multicolumn{2}{c}{Game Master} \\
    \hline

    \theutterance \stepcounter{utterance}  
    & & & \multicolumn{4}{p{0.6\linewidth}}{
        \cellcolor[rgb]{0.9,0.9,0.9}{
            \makecell[{{p{\linewidth}}}]{
                \texttt{\tiny{[P1$\langle$GM]}}
                \texttt{You are an expert in causal inference. The following question is not a typical commonsense query, but rather a meticulously designed question created by a professor specializing in causal inference, intended to assess the students' mastery of the course content.} \\
\texttt{Imagine a self{-}contained, hypothetical world with only the following conditions, and without any unmentioned factors or causal relationships: Jyka has a direct effect on yupt and kwox. Yupt has a direct effect on kwox. We know that jyka causes yupt. jyka or yupt causes kwox. Would an individual is kwox if not jyka instead of jyka?} \\
\texttt{Start your answer with 'Yes' or 'No', followed by additional reasoning or evidence to support your explanation.} \\
            }
        }
    }
     \\ \\

    \theutterance \stepcounter{utterance}  
    & \multicolumn{4}{p{0.6\linewidth}}{
        \cellcolor[rgb]{0.561,0.737,0.561}{
            \makecell[{{p{\linewidth}}}]{
                \texttt{\tiny{[P1$\rangle$GM]}}
                \texttt{No.} \\
\\ 
\texttt{To answer this question, we need to apply the principles of causal inference, specifically the concept of counterfactuals and the backdoor criterion.} \\
\\ 
\texttt{Given the conditions:} \\
\\ 
\texttt{1. Jyka has a direct effect on yupt and kwox.} \\
\texttt{2. Yupt has a direct effect on kwox.} \\
\texttt{3. Jyka causes yupt.} \\
\texttt{4. Jyka or yupt causes kwox.} \\
\\ 
\texttt{We want to know whether an individual would be in state kwox if not jyka instead of jyka. This is a counterfactual question, which asks what would have happened if a particular condition (in this case, jyka) had not occurred.} \\
\\ 
\texttt{To answer this question, we need to consider the causal relationships between the variables. Since jyka causes yupt, and yupt has a direct effect on kwox, we can say that jyka has an indirect effect on kwox through yupt.} \\
\\ 
\texttt{Now, let's consider the condition "jyka or yupt causes kwox." This means that either jyka or yupt (or both) can cause kwox. However, since jyka causes yupt, we can say that jyka is a sufficient cause of kwox through yupt.} \\
\\ 
\texttt{To determine whether an individual would be in state kwox if not jyka instead of jyka, we need to consider the backdoor criterion. The backdoor criterion states that if there is a path from the cause (jyka) to the effect (kwox) that is blocked by a collider (yupt), then the effect is not caused by the cause.} \\
\\ 
\texttt{In this case, the path from jyka to kwox is blocked by the collider yupt. Since jyka causes yupt, and yupt has a direct effect on kwox, the backdoor criterion is satisfied.} \\
\\ 
\texttt{Therefore, we can conclude that an individual would not be in state kwox if not jyka instead of jyka, because the effect of jyka on kwox is blocked by the collider yupt.} \\
            }
        }
    }
    & & \\ \\

    \theutterance \stepcounter{utterance}  
    & & & \multicolumn{2}{p{0.3\linewidth}}{
        \cellcolor[rgb]{0.9,0.9,0.9}{
            \makecell[{{p{\linewidth}}}]{
                \texttt{\tiny{[GM$|$GM]}}
                \texttt{No} \\
            }
        }
    }
    & & \\ \\

    \theutterance \stepcounter{utterance}  
    & & & \multicolumn{2}{p{0.3\linewidth}}{
        \cellcolor[rgb]{0.9,0.9,0.9}{
            \makecell[{{p{\linewidth}}}]{
                \texttt{\tiny{[GM$|$GM]}}
                \texttt{no} \\
            }
        }
    }
    & & \\ \\

    \theutterance \stepcounter{utterance}  
    & & & \multicolumn{2}{p{0.3\linewidth}}{
        \cellcolor[rgb]{0.9,0.9,0.9}{
            \makecell[{{p{\linewidth}}}]{
                \texttt{\tiny{[GM$|$GM]}}
                \texttt{game\_result = WIN} \\
            }
        }
    }
    & & \\ \\

\end{supertabular}
}

\end{document}
